% SPDX-FileCopyrightText: Copyright (c) 2026 NVIDIA CORPORATION & AFFILIATES. All rights reserved.
% SPDX-License-Identifier: Apache-2.0
%
% ROLE: what could make the results wrong inside AIS-timed replay, or keep them from transferring
%   to a live deployment, and what the campaign does about each. Organized as: simulator validity,
%   workload and objective validity, objective gaming, statistical validity, comparison validity,
%   deployment scope. Each item gives the threat, its likely direction where known, the
%   mitigation, and what stays open.
% STATUS: final, including the live GPU check (sec:live).

\section{Discussion, limitations and threats to validity}
\label{sec:threats}

Every claim this study makes is a claim about AIS-timed offline replay of one deployment. This
section lists what could make the headline wrong inside that setting, or keep it from transferring to
a live router, with what the campaign did about each and what stays open. \Cref{tab:threats}
summarizes; the subsections give the evidence.

\paragraph{What the results say.}
Within replay, the picture is consistent. Cache-aware routing matters a great deal on this deployment
(round-robin trails the default router by a segment-mean clipped log-ratio of $-0.56$), and the strong cache-aware
heuristics, once tuned, are close to each other. A conditional-logit policy over the router's own
signals, started from the default router, beats the branch's ported heuristics, each tuned with an
equal budget; the step that passes is the extended feature set, whose largest learned
term makes uncached prefill more expensive the busier a worker is. The margin is small, below the
pre-registered minimum detectable effect, and stable in sign under lag, timing, SLO-definition and
window perturbations. Three findings bound it. Most of the validation margin is reproduced by the
default cost plus a faithful LMetric term, which the baseline pool did not contain. The gain is concentrated on short and mid
prompts, paid for by long prompts. And an unconstrained search found a degenerate policy that only a
sign constraint removed. On GPUs, six test cells served live agree with the simulated ranking and
with the sign of M1-v2 minus \policy{ramjet}, but not with absolute goodput or with the size of the gaps
between policies (\cref{sec:live}).
% src: sec:results and sec:gaming (CR/facts/test_results.json, robustness.json,
%      fairness_in_class_heuristic.json; CR/audits/phase2-mid/gaming-concentration.md)

\begin{table}[tbp]
  \centering\footnotesize
  \caption{Threats to validity. ``Direction'' is the likely effect on the learned policy's measured
    advantage, where one can be argued.}
  \label{tab:threats}
  \begin{tabularx}{\linewidth}{@{}YYYY@{}}
    \toprule
    Threat & Likely direction & Mitigation & Status \\
    \midrule
    \ais{} timing differs from GPUs & unknown & live ranking check; timing $\times 0.8$, $\times 1.2$ &
      timing: sign holds; live: ranking and sign agree, absolute goodput and gap sizes do not \\
    Live check on few, chosen cells & inflates live gaps & registered analysis; validation only &
      six cells, chosen after the test pass; one run each \\
    Router state perfectly fresh in replay & favors policies that lean on exact state & lag
      10--\SI{200}{ms} on test & ranking holds \\
    No host, disk or global KV tier & excludes tier-dependent wins & scope statement & inherent \\
    AgentX recorded with another model, closed loop only & unknown & family reported separately &
      inherent \\
    SLOs and loads set with $\piref$ & unknown & SLO scale and definition re-scoring & sign holds;
      significance only at scales 1--3 \\
    Objective gaming by concentration & inflates & guards; audit; sign constraint & lost-requests
      worker removed; segregation reported \\
    12 independent test segments & low power at small gaps & pre-registered test; MDE & effect below
      MDE: sign only \\
    Selection and redesign on validation & inflates validation numbers & single frozen test pass;
      symmetric selection & disclosed \\
    Ports differ from papers & inflates & stated per port; in-class diagnostic & headline scoped to
      the ports as evaluated; \policy{lmetric} since fixed, not re-run \\
    Equal budget, unequal search & unknown & equal $\budget$ and restarts; disclosed starts & no
      baseline under-tuned \\
    \bottomrule
  \end{tabularx}
\end{table}

% =====================================================================================
\subsection{Simulator validity}
\label{sec:threats-sim}

\paragraph{Simulated timing only.}
Every latency in this paper comes from \ais{}'s forward-pass model; none was measured on GPUs, and
the study does not validate \ais{}'s absolute accuracy. What matters here is narrower: whether
replay ranks policies as GPUs would. That is not implied by aggregate accuracy. In a study of
learned routing for disaggregated serving, a simulator re-fitted to match measured goodput to a
mean absolute residual of 0.068 still did not reproduce the measured winner on any workload
\citep{tumkur2026calibrate}; that simulator was a hand-built discrete-event model, so how far the
lesson carries to \ais{}-timed engine replay is an open question. Controllers trained in emulation
have also lost to simple ones in the real world \citep{yan2020puffer}. Two features of this
campaign raise the stakes. The calibrated loads sit in the default router's transition region,
where performance-model errors grow fastest \citep{agrawal2024vidur}. And AgentX requests reach
128K tokens, where \ais{}'s accuracy has not been checked against hardware; the setup only showed
that replay's 120K-token prefill time equals an independent sum of \ais{}'s chunk estimates
(\SI{17815.54}{ms}, about 64\% model FLOPs utilization), an internal-consistency check. No AgentX
play fits a 32K context (caveat 10 of \cref{tab:caveats}), so this family cannot be moved to a
better-validated range. What exists: engine speed scaled by 0.8 and 1.2 on test, four simulator-model
counterfactuals that act on M1-v2's concentration mechanism (\cref{sec:res-robustness}), and one live
smoke run of the default router on one cell, which matched replay's windowed goodput within 0.25\%
while live prefill was slower and live decode faster than \ais{}'s, two errors that offset there
(\cref{sec:live-lane}). The finalist comparison on GPUs compares rankings rather than absolute goodput
(\cref{sec:live}). On six $\Nw = 4$ test cells the live ranking agrees with simulation and M1-v2 leads
\policy{ramjet} wherever simulation predicts it, but the two errors no longer offset everywhere: live
goodput ranges from 0.72 to 2.18 times the simulated value, differently for different policies in a
cell, and 18 of the 30 live deltas against $\piref$ depart from simulation beyond live noise, in both
directions. So replay's ranking held on hardware where it was checked, and its gap sizes did not.
% src: CR/facts/live_results.json wording, primary.departures (vs_default 18 of 30); CR/report/LIVE.md 15.2
% src: CR/facts/paper_draft.json derived_numbers.agentx_32k_fit; CR/facts/live_smoke.json verdict
%      (ratio 0.9976), discrepancies_live D4; CR/audits/final/refute-headline-3.md (simulator variants)
% src: CR/literature/LESSONS.md LR-14 (CtR §VI, MAE 0.068, "does not reproduce the measured winner";
%      Puffer §5.2; Vidur §7.2 and App. A.1) and "Risks the literature flags" 1, 11;
%      CR/facts/STATE.md audit(setup, determinism-cost-context, r0) ("120K TTFT 17,815.54 ms == independent
%      AIS chunk sum (~64% MFU, plausible)"); CR/CONTRACT.md A13

\paragraph{Perfectly fresh router state.}
Replay applies KV events to the router's index synchronously, so a policy always sees exactly which
blocks each worker holds; a live router sees those events late. Load counters (active prefill
tokens, decode blocks, active requests) come from the router's own request tracking in both
settings, but that their update timing matches is an assumption. Stale state matters most for a
deterministic argmax: routing to the least-loaded server on old information can perform much
worse than a randomized choice \citep{mitzenmacher1998old}, and greedy dispatch on delayed state
herds \citep{tahir2022meanfield}. The expected direction was that a policy leaning harder on exact
overlap and on small load differences, as the learned arms do, would lose more of its replay
advantage under lag than the baselines. Training never sees lag, by operator decision; the test stage
re-ran 17 finalists at $\lag \in \{10, 50, 200\}$\,ms on a separate build that delays KV events,
prefill completions and request completions by $\lag$ while keeping the router's own dispatch
bookkeeping immediate. The expectation was not borne out: M1-v2 kept first place at every lag with
its lead inside the nominal interval, and on validation at \SI{1000}{ms} the baselines lost more to
staleness than M1-v2 did (\cref{sec:res-robustness}). Two limits remain. The lag model is uniform and
deterministic, while live staleness varies with load and with router replicas; and a concentrated
policy is the kind most exposed to herding, which only the live runs can test.
% src: WT/lib/router-plugins/builtin/src/learned_choice/FEATURES.md "Information parity with the live router";
%      CR/CONTRACT.md A5.1, A11.5; CR/facts/lag_patch.json semantics, proof_lag0_identical (64/64, 404,704 rows),
%      proof_lag_changes_what_policies_see.goodput_level (min/max ratios), caveat;
%      CR/literature/LESSONS.md LR-14 (Mitzenmacher pdf p.2; Tahir pdf p.1-2)

\paragraph{One cache tier, and no global KV store.}
Replay models only the GPU tier. A live Dynamo router's accounting estimate blends host and disk
tiers into the effective overlap, so feature 1 (\code{overlap\_frac}) is the device overlap in
replay but a weighted blend live, a distribution the policy never saw. Replay also has no global KV
store, which changes which published routers can win. SMetric's advantage over the best baseline
holds ``except for setups without any global KV\$ store'', and in SMetric's own measurements on an
agentic workload without a global tier, LMetric's product score served 36\% less traffic within
SLO than a tuned linear cache-plus-load score \citep{wang2026smetric,zhang2026lmetric}. Session
migration is cheap only when a global tier holds the context, so gains of that kind cannot
appear here, and SMetric was not ported. The results speak to deployments whose KV cache lives on
the serving GPUs.
% src: FEATURES.md "Information parity" (host and disk tiers); CR/literature/LESSONS.md LR-04, LR-07
%      (SMetric pdf p.6: 920 vs 1,437 TPS; pdf p.10 quote), "Risks the literature flags" 15

\paragraph{Replay-specific inputs.}
Two replay artifacts could leak into a policy and were checked. Replay fills a request's expected
output length with the true one; \lc{} never reads that field, and a leakage audit confirmed it
(\cref{sec:integration}). Replay also gives every request of a closed-loop flat trace a synthetic
single-use session identifier that a live client would not send. For feature 6 this is harmless,
since a single-use session never has a previous worker, but it made a session-keyed home feature
of the optional set \code{v2} behave differently in open and closed loop until that feature was
re-keyed on the prompt prefix. The upstream fix is not in the simulator version the campaign pins.
% src: FEATURES.md "Excluded inputs", "Feature set v2" (hash_home paragraph); CR/facts/UPSTREAM_FOLLOWUPS.md #6
%      (fixed in aisimulate 1bf36cfb, not released into the pinned 0.13.0.dev202609300000000061)

% =====================================================================================
\subsection{Workload and objective validity}
\label{sec:threats-workload}

\paragraph{AgentX is another model's traffic.}
The 82 AgentX plays were recorded with Opus, not Qwen3-32B, and are replayed as Qwen3-32B traffic
at up to 128K tokens. Request lengths, tool-call gaps and the dependency structure between calls
are the recorded ones, so the family has realistic agent structure but not Qwen3-32B's own output
lengths or stopping behavior. Idle gaps longer than \SI{300}{s} are capped: 140 of the 2,666 gaps,
in 56 of the 82 plays, exceed it, and uncapped the plays would span \SI{399658}{s} of simulated
time instead of \SI{11942}{s}. Plays are recycled with disjoint hash ranges so that worker counts up
to 32 stay loaded, which by construction rules out KV reuse across copies. AgentX runs only in
closed-loop lanes (caveat 7 of \cref{tab:caveats}) and has no live load generator, so it is the one
family the live check cannot cover. It is also the family on which the pilot M1 trailed and on which
M1-v2 now leads on all four segments, so these choices bear directly on the headline.
% src: WT/notes/learned-routing/campaign/PLAN.md Workloads ("labeled as Opus-recorded agent traffic
%      projected onto Qwen3-32B"); CR/facts/DEVIATIONS.md build B2 (300 s cap: 140 of 2,666 gaps, 56 of 82 plays,
%      399,658 s vs 11,942 s); CR/CONTRACT.md A3 (recycling, disjoint hash ranges); CR/facts/gate.json
%      audit_flags_for_stage4 (M1 trails on all three AgentX segments)

\paragraph{Generated and closely related traces.}
Session traffic is generated: seeded multi-turn conversations whose turns extend the whole
context, with lognormal lengths and exponential think times measured from the previous turn's
completion. Replay's native session source could not express a growing conversation. The flat
traces, Mooncake and the two FAST25 traces, come from one provider, and a third public trace was
excluded as a relabeled copy of Mooncake. The FAST25 conversation trace shares the arrival and
length skeleton of two Mooncake windows, so it adds cells, not independent workloads.
% src: CR/facts/DEVIATIONS.md build B2 (synthetic sessions; workload families); CR/facts/UPSTREAM_FOLLOWUPS.md #5;
%      CR/facts/HEADLINE_TEST.json test_set.segment_note

\paragraph{Arrival burstiness is a modeling choice.}
Mooncake and FAST25 conversation timestamps are logged on a grid of about \SI{3}{s} (the generated
sessions on \SI{1}{s}), so a replay that keeps them puts whole bursts on one instant at any speed-up.
Each replicate instead re-draws every session's first arrival uniformly inside its logging slot.
The choice is large: at $\Nw = 2$ and light load, $\piref$'s 85th-percentile slowdown was 8.32 with
the logged bursts and 1.46 with spreading. Absolute goodput levels therefore depend on this
decision; paired comparisons on the same replicate are less exposed.
% src: CR/facts/DEVIATIONS.md calibration "arrival spread replicates (protocol crn-spread-v1)";
%      CR/facts/UPSTREAM_FOLLOWUPS.md #8

\paragraph{Open-loop arrivals do not react.}
In open-loop cells requests arrive on the trace's schedule whatever the latency, while a real
client slows down when served slowly. Trace-driven simulation is biased when the policy changes
the traffic it replays \citep{alomar2023causalsim}. Closed-loop and agentic-lane cells release
work causally and are reported separately.

\paragraph{SLOs and loads are defined by the default router.}
Each family's $(\Imax, \Smax)$ was fixed from $\piref$'s runs at $\Nw = 8$, and each cell's load
is where $\piref$'s good fraction falls to 0.95, 0.85 or 0.65. That puts every cell in the default
router's transition region, which is where routing matters, but the region is defined relative to
one baseline, and every ratio is normalized by the same baseline. The SLOs also differ widely across families ($\Smax$ from 1.16 for
AgentX to 3.47 for Mooncake), and the end-to-end slowdown clause does most of the work: the ITL
clause alone fails 0.6--4.5\% of $\piref$'s requests at L2. The objective is therefore close to
``finish near the uncontended latency'', and a different SLO could rank policies differently. The
re-scoring of the test records under four other definitions of a good request kept M1-v2 first
under each of them; under scaled SLOs it stayed first among the router-observable policies at every
scale, though at scales 2 and 3 the \ais{}-league M2-ais edged it overall. Significance holds only at
the calibrated scale and looser, and at a tighter SLO the session segments reverse
(\cref{sec:res-robustness}).
% src: CR/facts/test_results.json a14_sla_transfer.variants (rank_of_headline 1 for the four
%      definitions and scales 0.5-1.5, 2 at scales 2.0 and 3.0; vs_default_cell_mean m2ais 0.0643 vs
%      m1v2 0.0641 at 2.0, 0.0470 vs 0.0466 at 3.0); generated/robustness.tex (rank column)
% src: CR/facts/calibration.json sla.<family>.e2e_slowdown, itl_only_fail_frac_default_at_L2
%      (sessions 0.0063, AgentX 0.0178, Mooncake 0.0449); CR/CONTRACT.md A2.2, A14

\paragraph{Worker counts are compared at matched stress.}
Loads were matched to $\piref$'s knee at each $\Nw$, not to a fixed per-worker rate; across worker
counts the per-worker load differs by 1.2--1.9$\times$ (2.4$\times$ in AgentX prefill). The
held-out-$\Nw$ result therefore compares policies at equally stressed operating points of the
default, which is the intended question, but it does not show behavior at a matched per-worker
rate. At large $\Nw$ capacity pools and simple policies catch up, and at small $\Nw$ scarcity can
favor randomization, so per-$\Nw$ differences need that context \citep{vanderboor2022survey,tumkur2026calibrate}.
% src: CR/facts/STATE.md audit(calibration/independent-rerun r1) minor F3 ("knee-matched N-scaling only
%      (per-worker 1.2-1.9x across N; AgentX prefill 2.4x)"); CR/literature/LESSONS.md LR-12

% =====================================================================================
\subsection{Objective gaming and concentration}
\label{sec:threats-gaming}

Black-box search exploits whatever its objective leaves out \citep{lehman2020creativity}. Three
failure modes are documented for policies like this one. A binary per-request goodput rewards
giving up on requests that will miss their SLO anyway \citep{wang2024slo}. An evolution strategy
that tuned a routing cost for p95 TTFT alone set the queue weight to zero and sent all traffic to
one replica \citep{toniolo2026gorgo}. And learned load balancers can drift into reserving or
segregating servers in ways that break under workload shift \citep{mao2019park}. Here, the ITL
clause bounds the second failure, but only weakly, because a mean ITL averages out a long stall
\citep{wang2024slo}.

The audit found the second reading, sacrifice, for the first tuned M1, and the sign constraint
removed its mechanism (\cref{sec:gaming}). What remains is segregation by request size, which the
request-count objective rewards: M1-v2 packs short cached requests onto a hot worker and gives up some
long-prompt goodput. Whether that is acceptable depends on what the deployment values. Under the
operator's objective it is a legitimate win; under a size-aware objective it is partly a transfer,
as the mixed token-weighted goodput shows (\cref{sec:a19}). Concentration also interacts with stale
state: a policy that funnels traffic to a hot worker is the kind most exposed to herding, and only
uniform, deterministic lag was tested (\cref{sec:threats-sim}).
% src: CR/audits/phase2-mid/gaming-concentration.md F1, F2; CR/audits/phase2-mid/fix.md "Re-audit result";
%      CR/CONTRACT.md A19 (operator decision)

% =====================================================================================
\subsection{Statistical validity}
\label{sec:threats-stats}

\paragraph{Few independent units.}
The headline rests on 12 test segments, the unit of independence, and per-family claims rest on 2--4
and are descriptive only \citep{demsar2006}. The segment structure is uneven: Mooncake and FAST25
conversation contribute 28 of the 60 test cells but 2 of the 12 segments, AgentX 14 cells and 4
segments (\cref{tab:headline}). Workload replicates and duplicated trace copies lengthen and
re-draw a segment; they do not add independent segments.
% src: CR/facts/HEADLINE_TEST.json test_set; CR/cells/test.jsonl (segment counts); CR/CONTRACT.md A4

\paragraph{Effects near the detection limit.}
The headline effect, $+0.0316$ per segment, is below the minimum detectable effect of 0.038 that the
pilot fixed, and below the 0.060 spread across timing perturbations. The test still rejects, because
10 of 12 segments favor M1-v2 and the two that do not are ties at ceiling; but a significant small
effect is not a measured size, so only the sign is claimed. The MDE is also measured in units of
(ratio $-\,1$), while the headline statistic is a clipped log-ratio; at these sizes the two agree to
within a few percent of the effect ($\log 1.038 = 0.0373$).
% src: CR/facts/test_results.json headline (segment_mean_vs_mde -0.0066); CR/facts/robustness.json
%      lr14_summary (spread 0.060); CR/facts/pilot.json mde

\paragraph{Selection on validation.}
Many configurations are compared on validation, and selection overfits what it selects on
\citep{cawley2010overfitting}. Disjoint uses of data contain this (\cref{sec:hygiene}): each
policy's configuration is chosen on validation replicates $\krep = 0,1,2$; the best baseline, and
by the same rule the headline learned arm, are chosen on fresh validation replicates
$\krep = 3,\dots,10$; and the test set is evaluated once. The cross-policy step reuses the same
fresh replicates for both choices, which keeps the comparison symmetric but makes both selected
validation scores optimistic; only the test numbers are clean. Every learned variant scored on validation is counted and reported
\citep{dodge2019show}, and the 7 test cells at $\Nw = 6$, a worker count validation also used, are
labeled selection-exposed. With 6 validation segments, selection itself is noisy: the choice of
\policy{ramjet} over \policy{llm-d-precise-prefix} flipped between the cell and the segment mean, and
on test \policy{llm-d-precise-prefix} was the stronger baseline. The learned side was also redesigned
on validation evidence while the baseline side was not (\cref{sec:res-fairness}); the single test pass
keeps the headline's $p$-value valid for the selected arm, but the originally planned M1 and M2 would
not have passed.
% src: CR/facts/gate.json full_plan.objective_cells_repeats.selection, selection_hygiene; CR/CONTRACT.md A11.3, A17.2;
%      CR/facts/HEADLINE_TEST.json counting

% =====================================================================================
\subsection{Comparison validity and fairness}
\label{sec:threats-fairness}

\paragraph{Ports, not papers.}
Baselines run as implemented on the campaign branch; by operator decision there are no
paper-faithful variants, although different implementations of one method can perform
materially differently \citep{henderson2018matters}. The known gaps and whether
tuning can close them:
\begin{itemize}
  \item \policy{lmetric} scores uncached prompt tokens times active requests plus one, without the
    queued-prefill term the paper credits for its advantage \citep{zhang2026lmetric}. Tuning cannot
    add it: the port's tunable parameters are its hot-spot detector's, and at the largest
    allowed step they change 0.06\% of decisions. The queued-prefill product is representable in
    the learned model's feature set \code{v2}, which the headline arm uses; this asymmetry is why
    the headline is scoped to the branch's ports: the default cost plus faithful LMetric carries most
    of the validation margin, and faithful LMetric alone, untuned, beats every tuned baseline on
    validation (\cref{sec:res-fairness}). This describes the port as evaluated. After the campaign
    froze its results, the port in ai-dynamo/dynamo\#15450 was fixed to include the queued-prefill
    term, so it now matches the paper's definition; every \policy{lmetric} number here is from the
    port before the fix, the headline's scope refers to the ports as evaluated, and no equal-budget
    comparison with the fixed port has been run.
  \item \policy{dualmap} ships a pending-prefill budget of \num{65536} tokens and
    \policy{llm-d-optimized-baseline} a peak prefill rate of \num{15928} tokens/s, both from other
    deployments \citep{yuan2026dualmap}. Tuning reaches both: the budget is a free parameter, and
    the rate is pinned only because its product with the free TTFT-penalty parameter is what
    determines decisions.
\end{itemize}
A brief sanity check found that every port's decisions respond to its inputs in the documented
direction; two of them (\policy{dualmap}, \policy{chwbl}) lost to the default router at shipped
defaults on the pilot's train subset, plausibly because hash-pinned policies suffer when a few hot
prefixes carry most of the traffic.
% src: CR/CONTRACT.md A2.5; CR/literature/LESSONS.md LR-04 ([code] lmetric.rs:151-157, dualmap.rs:57,
%      optimized_baseline.rs:79); CR/facts/phase2_prep.json spaces (lmetric free: class_prefix_blocks,
%      window_requests; flip_at_sigma0 0.00058; dualmap free incl. pending_prefill_token_budget;
%      llmd_optimized_baseline free incl. max_ttft_penalty_ms), decisions (gauge pins);
%      CR/facts/pilot.json heuristic_screen_train_subset.sanity_note; FEATURES.md v2 representability

\paragraph{Equal budget, unequal search.}
Every tuned policy gets $\budget = 400$ evaluations per restart and $\restarts = 3$ restarts, but
the policies differ in what that budget can buy. The search spaces range from 2 free parameters
(\policy{lmetric}, \policy{chwbl}) to 7 (M1) and 22 (M1-v2). Each space's initial step size was set so that one
step changes about 15--20\% of routing decisions; five baselines cannot reach that even at the
largest allowed step, where they change 0.06\% (\policy{lmetric}), 1.8\% (\policy{dualmap}), 3.4\%
(\policy{llm-d-optimized-baseline}), 8.3\% (two-tier) and 13.3\% (\policy{ramjet}) of decisions.
For those policies tuning cannot move far from the shipped behavior; that is a property of their
parameterization, and a learned model with a richer one is expected to gain more from the same
budget. The equal-budget rule fixes what each policy may spend, not how far each can get. The
fairness audit's validation-oracle probes found no baseline that more search would have improved
(\cref{sec:budget}).
% src: CR/facts/phase2_prep.json spaces.<name>.{dimension, sigma0, flip_calibration.flip_at_sigma0_interpolated,
%      target_reached_by_0.3}; CR/facts/gate.json full_plan.budget

\paragraph{Initialization.}
Baselines restart from their shipped defaults. The learned arms' second and third restarts start
from a behavior clone of \policy{llm-d-precise-prefix} and from a third informed point, for M1-v2 the
faithful LMetric point. These starts use train-split decisions, not extra objective evaluations, so
the budget is unchanged, but they hand the learned arms information about strong heuristics. Both
selected configurations, M1's and M1-v2's, came from the default-router start s1, and M1 restricted
to its default start is the same policy as M1 (\cref{sec:res-ladder}). The faithful LMetric start is
nevertheless strong on its own, which is the scope issue of \cref{sec:res-fairness}.
% src: CR/CONTRACT.md A11.1, A12; CR/facts/test_results.json secondary.m1_A12_note;
%      CR/facts/finalists.json headline.learned (m1v2-s1)

\paragraph{Excluded and separated baselines.}
ThunderAgent's routing half needs a request classifier that replay cannot run and is not
evaluated \citep{kang2026thunderagent}. Policies that read \ais{} estimates at decision time,
including llm-d's optimized baseline with modeled TTFT, compete in a separate, secondary league
against the same normalizer (\cref{sec:res-ais}). SMetric and faithful LMetric were not tuned as
baselines; an equal-budget comparison with them on fresh data is the open follow-up of
\cref{sec:res-fairness}.
% src: CR/CONTRACT.md A16; WT/notes/learned-routing/campaign/PLAN.md Baselines

% =====================================================================================
\subsection{Deployment scope}
\label{sec:threats-scope}

The study covers one model (Qwen3-32B), one engine (vLLM 0.24.0), one GPU type (H100 SXM),
tensor parallelism 2, aggregated serving, identical workers, and 2 to 32 workers. Disaggregated
prefill and decode, other models and hardware, and heterogeneous workers are out of scope. The
coefficients are specific to this deployment; moving to another means re-tuning, and the cost of
doing so (\cref{tab:cost}) is part of what a team would weigh.

The headline policy reads only inputs the live router already computes, so its YAML runs in the
live router unchanged; whether those inputs carry the same values and timing live is what the lag
check and the live run test. It decides by argmax ($\temp = 0$) and leaves the router's queue
threshold unset, so it never holds a request back; a randomized variant and a jointly tuned queue
threshold were tested and added nothing. Shipping it needs four settings to match replay: the
router's structural knobs (\code{router\_track\_prefill\_tokens}, \code{router\_assume\_kv\_reuse},
\code{router\_track\_active\_blocks}) at their defaults, since they change what the features mean
(caveat 3 of \cref{tab:caveats}); host session affinity off, or accepted as a replacement for
$\theta_6$, since the policy asks the host to treat an affinity target as exclusive; no \code{seed},
so that router replicas do not break ties in lockstep; and monitoring of the concentration guards
of \cref{sec:gaming}, since the policy's expected behavior includes a hot worker at about twice the
mean in-flight load. A live router's overlap estimate also blends host and disk tiers, which replay
does not have. The simulator-informed league would need \ais{} estimates, or a closed form
calibrated offline, inside the live router.
% src: CR/report/REPORT.md section 12 (porting notes); CR/audits/build/leakage-features-r1.md F1, F4
% src: CR/CONTRACT.md "Deployment (operator-fixed)", A5.2, A16.1 (info parity of AIS signals);
%      FEATURES.md "Information parity"; CR/facts/gate.json full_plan.tiers.C_conditional
